\clearpage\fancyhead[L]{\small COMET.}\hypertarget{native-338-lower}{}\label{native-338-lower}
\begin{center}\textbf{COMET.}\par\vspace{4pt}\hrule\vspace{4pt}\fontsize{10}{12}\selectfont Astronomisches Jahrbuch für 1814. Seite 254\,..\,256. Berlin 1811.\par\vspace{4pt}\hrule\vspace{4pt}\end{center}

\newcommand{\ObservedNumeratorMark}{\hbox{\vrule width0pt height2.291667bp depth0pt\special{pdf:bcontent}\special{pdf:code q 0 g 1.208333 0.833333 m 1.083333 0.833333 l 1.041667 1.000000 l 0.708333 1.000000 l 0.708333 0.875000 l 0.541667 0.833333 l 0.541667 0.708333 l 0.375000 0.666667 l 0.333333 0.500000 l 0.000000 0.500000 l 0.000000 0.000000 l 0.500000 0.000000 l 0.541667 0.166667 l 0.875000 0.166667 l 0.875000 0.291667 l 1.041667 0.333333 l 1.041667 0.500000 l 1.208333 0.541667 l h 0.000000 2.083333 m 0.000000 1.958333 l 0.125000 1.958333 l 0.166667 1.791667 l 1.041667 1.791667 l 1.083333 1.583333 l 1.375000 1.583333 l 1.375000 2.083333 l 1.208333 2.125000 l 1.208333 2.291667 l 0.166667 2.291667 l 0.166667 2.125000 l h f* Q}\special{pdf:econtent}\kern1.375000bp}}

\begin{flushright}Göttingen den 10. September 1811.\end{flushright}
\par\indent Den Cometen haben wir hier zuerst den 22. August in der Abenddämmerung gesehen, allein die von der Stadt begrenzte Aussicht unserer Sternwarte erlaubte erst im September Beobachtungen zu machen. Ich habe mich dazu der Abstände von Fixsternen bedient, die sich mit dem Sextanten vom Stativ sehr gut beobachten lassen. Alle bisher gemachten Beobachtungen haben noch nicht reducirt werden können; ich kann heute blos erst folgende drei mittheilen, denen Abstände von $\alpha$ Lyrae und $\alpha$ Aurigae zum Grunde liegen.
\Needspace{9\baselineskip}
\begin{center}\fontsize{10}{12}\selectfont\setlength{\tabcolsep}{4pt}\renewcommand{\arraystretch}{1.15}\begin{tabular}{r|c|c|c}
1811 & Mittlere Zeit & Gerade Aufst. & Nördl. Abw.\\
\hline
Septbr. 4 & $8^{\mathrm u}\,28^{\mathrm m}\,47^{\mathrm s}$ & $158^\circ\,25'\,24''$ & $39^\circ\,18'\,2''$\\
6 & $8\;48\;38$ & $160\;23\;16$ & $40\;14\;16$\\
7 & $8\;57\;6$ & $161\;26\;10$ & $40\;41\;54$\\
\end{tabular}\end{center}

\par\indent Ich habe hienach und nach neuen Beobachtungen des Hrn. v. \textsc{Zach} meine vorläufigen parabolischen Elemente, die im 130. Stück uns. gel. Anz. [S. 327 d. B.] abgedruckt sind, folgendermassen verbessert.
\Needspace{11\baselineskip}
\begin{center}\fontsize{10}{12}\selectfont\setlength{\tabcolsep}{4pt}\renewcommand{\arraystretch}{1.15}\begin{tabular}{l}
Durchgang durch das Perihelium 1811 Sept. 12\\
\quad $5^{\mathrm u}\,21^{\mathrm m}\,15^{\mathrm s}$ Mittl. Zeit in Göttingen\\
\end{tabular}\end{center}

\begin{center}\fontsize{10}{12}\selectfont\setlength{\tabcolsep}{4pt}\renewcommand{\arraystretch}{1.15}\begin{tabular}{llr}
Logarithm. des kleinsten Abstandes & .\;.\;.\;.\;.\;.\;.\;.\;. & $0.017060$\\
Länge des Perihels & .\;.\;.\;.\;.\;.\;.\;.\;.\;.\;.\;.\;.\;.\;. & $75^\circ\,17'\,34''$\\
Länge des aufsteigenden Knoten & .\;.\;.\;.\;.\;.\;.\;.\;.\;. & $140\;24\;13$\\
Neigung der Bahn & .\;.\;.\;.\;.\;.\;.\;.\;.\;.\;.\;.\;.\;.\;. & $73\;7\;17$\\
\end{tabular}\end{center}
\noindent Bewegung rückläufig.

\Needspace{19\baselineskip}\SourceLeaf{339}\par\indent Ich habe hienach den geocentrischen Lauf auf 4 Monate vorausberechnet, wie folgt:

\begin{center}\fontsize{10}{12}\selectfont\setlength{\tabcolsep}{4pt}\renewcommand{\arraystretch}{1.15}\begin{tabular}{rr|c|c|c|c}
\multicolumn{2}{c|}{Göttinger Zeit} & \shortstack{Gerade\\Aufsteig.} & \shortstack{Nördliche\\Abweich.} & \shortstack{Log. der Dist.\\von der Erde} & \shortstack{Licht-\\stärke.}\\
\hline
1811 Sept. 12 & $5^{\mathrm u}$ & $166^\circ\,54'$ & $42^\circ\,58'$ & $0.2079$ & $0.355$\\
17 & $5$ & $173\;46$ & $45\;15$ & $0.1836$ & $0.394$\\
22 & $6$ & $182\;2$ & $47\;17$ & $0.1594$ & $0.432$\\
Octob. 2 & $7$ & $203\;24$ & $49\;34$ & $0.1160$ & $0.490$\\
12 & $8$ & $228\;45$ & $47\;15$ & $0.0911$ & $0.492$\\
22 & $9$ & $251\;23$ & $39\;56$ & $0.0963$ & $0.421$\\
Nov. 1 & $10$ & $267\;58$ & $30\;32$ & $0.1312$ & $0.310$\\
11 & $11$ & $279\;38$ & $22\;18$ & $0.1851$ & $0.208$\\
21 & $12$ & $288\;13$ & $14\;55$ & $0.2415$ & $0.138$\\
Dec. 1 & $13$ & $294\;54$ & $9\;42$ & $0.2976$ & $0.093$\\
11 & $14$ & $300\;24$ & $5\;49$ & $0.3487$ & $0.064$\\
21 & $15$ & $303\;22$ & $3\;0$ & $0.3848$ & $0.048$\\
\end{tabular}\end{center}

\par\indent Von der Mondfinsterniss am 2. Sept. kam hier wegen sehr ungünstigen Wetters nicht viel zu Gesicht, aber die während derselben eingetretene Bedeckung von $\lambda$ \emph{Aquarii} (im Jahrbuch etwa 1 Stunde zu früh angegeben) wurde sehr gut beobachtet. Ein- und Austritt geschahen am verfinsterten, aber doch noch sichtbaren Mondrande, und letzterer konnte um so besser beobachtet werden, da einer meiner Zuhörer diese Bedeckung voraus berechnet hatte, und wir also darauf vorbereitet waren.

\par\indent Eintritt $\lambda$ Aquarii Mittl. Zeit $10^{\mathrm u}\,16^{\mathrm m}\,17^{\mathrm s}4$. Austritt $10^{\mathrm u}\,51^{\mathrm m}\,40^{\mathrm s}5$.

\par\indent Noch füge ich Ihrem Wunsche zufolge einen kleinen Zusatz zu meiner Theoria motus corporum coelestium bei.
\Needspace{16\baselineskip}\begin{center}\emph{Zusatz zu Art. 90 und 100 der Theoria motus corporum coelestium, vom Verfasser.}\end{center}
\par\indent Zur Auflösung der wichtigen Aufgabe, aus zweien Radiis vectoribus und dem eingeschlossenen Winkel die elliptischen oder hyperbolischen Elemente zu bestimmen, habe ich mich mit grossem Vortheil einer Hülfsgrösse $\xi$ bei der Ellipse, $\zeta$ bei der Hyperbel bedient, für welche ich jenem Werke eine Tafel angehängt habe. Berechnet ist diese Tafel nach einem dort angeführten continuirten Bruche, dessen vollständige Ableitung aber dort nicht gegeben ist, und zu dessen theoretischer Entwickelung, die mit andern Untersuchungen zusammenhängt, ich bisher noch nicht Gelegenheit gefunden habe. Es wird daher manchem lieb sein, hier einen andern Weg angezeigt zu finden, auf welchem man jene Hülfsgrösse eben so bequem hätte berechnen können.

\par\indent Wir haben (Art. 90)
\[
\xi=x-\frac{5}{6}+\frac{10}{9X}=\frac{xX-\dfrac{5}{6}X+\dfrac{10}{9}}{X}
\]
\par\indent Der Zähler dieses Bruchs verwandelt sich leicht, wenn man für $X$ die dort gegebene Reihe substituirt, in

\[
\begin{aligned}
&\frac{8}{105}xx\bigl(1+\frac{2\mathord{.}8}{9}x+\frac{3\mathord{.}8\mathord{.}10}{9\mathord{.}11}xx\\
&\qquad+\frac{4\mathord{.}8\mathord{.}10\mathord{.}12}{9\mathord{.}11\mathord{.}13}x^3+\frac{5\mathord{.}8\mathord{.}10\mathord{.}12\mathord{.}14}{9\mathord{.}11\mathord{.}13\mathord{.}15}x^4+\text{etc.}\bigr)
\end{aligned}
\]
\par\indent Setzt man also die Reihe
\[
1+\frac{2\mathord{.}8}{9}x+\frac{3\mathord{.}8\mathord{.}10}{9\mathord{.}11}xx+\text{etc.}=A
\]
\SourceLeaf{340}\noindent so wird
\[
xX-\frac{5}{6}X+\frac{10}{9}=\frac{8}{105}Axx
\]
\[
X=\frac{\dfrac{4}{3}\bigl(1-\dfrac{12}{175}Axx\bigr)}{1-\dfrac{6}{5}x}
\]
\[
\xi=\frac{\dfrac{\mathord{\raisebox{7bp}{\ObservedNumeratorMark}\,}2}{35}Axx\bigl(1-\dfrac{6}{5}x\bigr)}{1-\dfrac{12}{175}Axx}
\]
\noindent nach welcher Formel man $\xi$ immer bequem und sicher berechnen kann. Für $\zeta$ braucht man nur $-z$ statt $x$ zu setzen.

\par\indent Ich bemerke nur noch, dass man $A$ noch bequemer nach folgender Formel berechnen kann

\[
\begin{aligned}
A={}&(1-x)^{-\frac{3}{2}}\bigl(1+\frac{1\mathord{.}5}{2\mathord{.}9}x+\frac{1\mathord{.}3\mathord{.}5\mathord{.}7}{2\mathord{.}4\mathord{.}9\mathord{.}11}xx\\
&\qquad+\frac{1\mathord{.}3\mathord{.}5\mathord{.}5\mathord{.}7\mathord{.}9}{2\mathord{.}4\mathord{.}6\mathord{.}9\mathord{.}11\mathord{.}13}x^3+\text{etc.}\bigr)
\end{aligned}
\]
\noindent allein die Ableitung dieser Reihe aus der vorigen beruht auf Gründen, die hier nicht ausgeführt werden können.
